%!TEX root = ../dokumentation.tex

\chapter{Discussion}
\label{ch:discussion}

\section{Answering the Research Questions}
\label{sec:rqs}

\textbf{RQ1: How does a classical RAG pipeline compare to an agent-based approach in terms of latency, cost, and retrieval performance in enterprise environments?}

The RAG pipeline outperforms the agent-based system on all three dimensions. On retrieval performance, RAG achieves a context precision of 0.546 on factual queries and 0.730 on synthesis queries, compared to the agent's 0.338 and 0.315 respectively; it also produces higher answer relevancy on both query types. On latency, the RAG pipeline completes queries in a mean of 2.51 seconds, against the agent's 16.33 seconds — a factor of 6.5. On cost, the agent consumes 15.7 times as many input tokens as RAG over the 90-query evaluation, translating to an estimated cost ratio of 11.2 to~1. These are not marginal differences: the agent's disadvantage on each dimension is large enough to dominate the decision in any deployment where latency budgets or token costs are constrained. RAG's advantage on retrieval performance is partly attributable to the agent's 41\% query failure rate, but even absent those failures, the structural differences in latency and cost reflect the multi-step architecture of the ReAct loop and would persist with a higher iteration ceiling.

\textbf{RQ2: To what extent does an agent-based system improve source selection and answer quality compared to RAG?}

The agent does not improve answer quality relative to RAG on any measured metric. However, it does provide a meaningful routing capability that RAG lacks by design: 73.3\% of queries are directed to the correct source type without any explicit routing rules, using only the natural-language tool descriptions as the selection mechanism. This is a genuine architectural advantage — RAG cannot distinguish which source type produced a retrieved chunk, and cannot selectively query one source to the exclusion of others. Whether this capability translates to better answers depends on the query mix. For factual queries against a single source, the agent's routing adds latency without benefit, because the unified RAG index would have retrieved the same chunk. For cross-source synthesis queries, the agent's ability to query each source independently and synthesise across observations is structurally superior, though the evaluation results on synthesis queries do not confirm a quality advantage in practice, partly because 41\% of queries — concentrated in the synthesis category — returned empty answers.

\textbf{RQ3: What are the differences in implementation complexity and maintainability between both approaches?}

The rubric assessment in Section~\ref{sec:complexity} assigns RAG an unweighted mean score of 4.0 out of 5 across the five complexity criteria, against the agent's 2.2. The most substantial differences are in debuggability and operational overhead. The RAG pipeline's linear three-stage data flow — retrieve, rerank, generate — makes any failure traceable to a single stage through the structured log. The agent's multi-step loop surfaces failures through non-deterministic tool selection, context accumulation, and the max-iterations boundary, each requiring different diagnostic procedures. In code volume, the agent implementation required 563 lines against RAG's 377 — a 49\% increase — primarily driven by the tool layer and the adaptations to the pipeline required for multi-step trace logging. Extensibility likewise favours RAG: adding a new data source requires two steps, compared to five for the agent, with the additional steps in the agent case spanning four files and necessitating revalidation of routing behaviour across all existing source types.

\textbf{RQ4: Under which conditions is each approach preferable in practice?}

The RAG pipeline is preferable in deployments where latency and cost are primary constraints, where the query mix consists largely of factual lookups against a single knowledge source, where the document corpus changes frequently (requiring fast index updates), and where the engineering team has limited experience with agentic frameworks. It is also preferable when the corpus is reasonably homogeneous — if all source types use the same document structure and terminology, a unified index retrieves across them effectively without explicit routing. The agent-based approach offers advantages when the query mix is dominated by multi-source synthesis questions that require combining information from structurally different knowledge systems, when the cost and latency premium is acceptable relative to the quality of synthesis it can — when it completes — produce, and when routing transparency is valued: the agent's \texttt{sources\_used} field in the log provides an auditable record of which source each answer drew from, which RAG cannot offer. In enterprise environments with strict token budgets or latency SLAs, the agent's 11.2$\times$ cost premium and 6.5$\times$ latency premium represent barriers that are difficult to justify without clear evidence of proportional quality gains — evidence that this evaluation did not produce.


\section{Trade-Off Analysis}
\label{sec:tradeoffs}

Table~\ref{tab:tradeoff} summarises the comparative position of both systems across the five evaluation dimensions. Ratings are qualitative assessments anchored to the specific numbers reported in Chapter~\ref{ch:results}.

\begin{table}[ht]
\centering
\caption{System-level trade-off comparison across five evaluation dimensions.}
\label{tab:tradeoff}
\begin{tabular}{lll}
\hline
\textbf{Dimension} & \textbf{RAG} & \textbf{Agent} \\
\hline
Retrieval quality      & High (CP: 0.55--0.73) & Moderate (CP: 0.31--0.34; 41\% failure rate) \\
Answer quality         & High (AR: 0.68--0.79) & Moderate (AR: 0.55--0.66; 41\% failure rate) \\
Latency                & Low (mean 2.5s)        & High (mean 16.3s; stdev 8.2s) \\
Token cost             & Low (\$1.18/90 queries) & High (\$13.24/90 queries; 11.2$\times$) \\
Implementation effort  & Moderate (377 LoC)    & High (563 LoC; 5 files; routing revalidation) \\
\hline
\end{tabular}
\end{table}

The trade-off is not symmetric. RAG's advantages — lower cost, lower latency, simpler implementation, and comparable or superior quality on the metrics available — obtain unconditionally in this evaluation. The agent's single structural advantage — explicit, auditable source routing — is a capability difference rather than a quality difference: it tells the user which source was consulted, but in this evaluation it did not produce better answers. The agent's routing accuracy of 73.3\% means that roughly one in four queries directed to the wrong source, which would typically produce a worse answer than the unified-index RAG approach that retrieves the best-matching chunks regardless of source.

Three deployment scenarios are identifiable where the agent's advantages would weigh more heavily. First, in a corpus where source types have genuinely non-overlapping terminology — for example, formal policy documents alongside engineering issue trackers — the agent's ability to direct a query to the correct source type may produce retrieval results that would not surface from a unified index trained on all documents simultaneously. Second, in high-stakes, low-volume deployments — such as regulatory compliance review — where answer completeness and source attribution matter more than throughput or cost, the 11.2$\times$ cost difference becomes less significant relative to the value of a confirmed multi-source answer. Third, in workflows where downstream automation consumes the \texttt{sources\_used} field for audit or escalation routing, the agent's explicit source trace has value independent of answer quality. None of these scenarios were confirmed by the evaluation, but the architectural properties of the agent make them plausible candidates for future investigation.


\section{Limitations}
\label{sec:limitations}

\textbf{Proxy data sources.} The corpus consists of GitHub Issues, Markdown documentation pages, and CNCF PDF whitepapers — open, public datasets chosen to simulate the Jira, Confluence, and SharePoint systems found in enterprise environments. Real enterprise data differs in several ways that may affect generalisability: proprietary terminology reduces cross-source retrieval overlap; access control policies fragment documents across systems in ways not captured by a static corpus; and document quality is more variable, including partially complete drafts, archived pages, and duplicate entries. Results obtained on the proxy corpus may therefore be more favourable to retrieval-based approaches than results on real enterprise data, where noise and formatting inconsistency would test chunking and embedding robustness more severely.

\textbf{Synthetic evaluation query set.} The 90 evaluation queries were generated by prompting \texttt{gpt-oss-120b} from sampled corpus excerpts and manually reviewed for quality. Synthetic query sets are vulnerable to a systematic bias: the generating model may produce questions that match the phrasing and structure of passages in the corpus, making retrieval artificially precise. The manual review process applied the criterion that the question wording must not quote or paraphrase the answer, which partially mitigates this risk, but cannot eliminate it entirely. A query set generated by domain experts or drawn from real user logs would provide a stronger evaluation baseline.

\textbf{Agent faithfulness unavailable.} As noted in Section~\ref{sec:robustness}, the RAGAS faithfulness metric could not be computed for any agent query because the LLM judge timed out on the agent's longer responses. Faithfulness — the degree to which answer claims are directly supported by retrieved context — is the metric most directly associated with hallucination risk. Its absence for the agent is a significant gap in the comparative analysis, particularly because the agent synthesises information across multiple tool observations and might be expected to introduce more unsupported claims than RAG's single-context generation.

\textbf{Agent max-iterations failures.} The 37-query failure rate (41\%) at \texttt{MAX\_ITERATIONS}=15 distorts all agent quality metrics, since empty-answer queries contribute a score of zero to BERTScore and ROUGE-L. A production system would address this through a query classifier that routes clearly unanswerable or scope-exceeding queries to a simpler fallback, or by implementing a recovery mechanism that extracts a partial answer from intermediate trace steps before the iteration ceiling is reached. Neither mechanism was implemented in the evaluation system, which means the agent's quality figures understate what the agent could achieve on the queries it successfully completes.

\textbf{Single-evaluator rubric.} The complexity and maintainability rubric scores in Section~\ref{sec:complexity} were assigned by the author based on the implementation experience described in Chapter~\ref{ch:implementation}. The scores are justified by objective proxies — lines of code, parameter count, extensibility step count — but they are not peer-reviewed or validated against a second independent implementation. The rubric mean of 4.0 versus 2.2 should be treated as a directional finding rather than a precise measurement.

\textbf{Single LLM evaluated.} Both systems use \texttt{gpt-oss-120b} as the generating model. This model is a reasoning model with internal chain-of-thought, which is atypical compared to the instruction-tuned models commonly evaluated in the RAG literature. The reasoning model's thoroughness — issuing up to nine tool calls per query — is the primary cause of the agent's high token consumption and max-iterations failures. A non-reasoning instruction-tuned model might produce a very different agent quality and latency profile. The results are therefore specific to the \texttt{gpt-oss-120b} model and may not generalise to other LLMs.


\section{Connection to Prior Work}
\label{sec:priorwork}

This evaluation contributes to a growing body of empirical comparisons between retrieval-augmented and agentic approaches to knowledge retrieval. Within the RAG taxonomy developed by \citeasnoun{gao2023ragsurvey}, the classical pipeline implemented here corresponds to Advanced RAG: it applies pre-retrieval chunking with overlap, embedding-based retrieval, and a cross-encoder reranking stage. The agent-based system corresponds to Agentic RAG — a sub-category of Modular RAG in which retrieval is dynamically routed across multiple specialised indexes rather than executed once against a unified index. The findings reported here are consistent with Gao et al.'s analysis of the trade-offs between these paradigms: Advanced RAG is computationally efficient and reliable but lacks the adaptive retrieval capability of agent-based approaches, while agentic variants offer greater flexibility at substantially higher cost.

The ReAct agent implementation follows the framework proposed by \citeasnoun{yao2023react} closely. Yao et al. demonstrate that interleaving reasoning traces with tool actions reduces error propagation and enables recovery from intermediate failures — properties confirmed in this evaluation through the agent's use of the \texttt{agent\_trace} to identify and discard hallucinated tool calls before computing evaluation metrics. However, the present results do not reproduce the quality advantages that Yao et al. report for ReAct over chain-of-thought prompting in their benchmark tasks. The difference is likely attributable to the nature of the retrieval task: Yao et al. evaluate question answering over Wikipedia, where a single search step typically suffices, whereas this evaluation includes synthesis queries that require combining evidence from structurally different document types — a harder retrieval problem that the agent does not reliably solve within the iteration ceiling.

The findings are most directly comparable to those of \citeasnoun{ferrazzi2026agenticrag}, who evaluate agentic RAG against classical RAG across enterprise-relevant corpora and report a similar pattern: agentic systems provide no consistent quality improvement over well-tuned classical RAG, while incurring substantially higher latency and cost. The present thesis extends Ferrazzi et al.'s analysis in two directions. First, the corpus is deliberately heterogeneous across structurally different document types — issues, documentation, and whitepapers — which simulates the multi-system nature of real enterprise knowledge environments more closely than a corpus drawn from a single source type. Second, implementation complexity is treated as an explicit evaluation dimension rather than a background consideration, producing the five-criterion rubric and its quantified scores. The rubric's finding — a complexity mean of 4.0 for RAG against 2.2 for the agent — provides a structured complement to the quality and performance results, and supports the same practical conclusion that Ferrazzi et al. reach through a different methodological path. The \citeasnoun{agenticrag2025enterprise} framework further contextualises these findings within enterprise deployment requirements, noting that agentic systems require additional governance mechanisms — source attribution, failure logging, iteration monitoring — that the structured logging implemented in this thesis begins to address through the \texttt{agent\_trace} and \texttt{sources\_used} fields.

Taken together, these findings suggest that the cost of the agentic architecture is not offset by quality gains in the setting evaluated here, and that the conditions under which the agent's structural advantages translate to measurable user benefit — heterogeneous terminology, mandatory source attribution, low-volume high-stakes queries — remain to be validated empirically.
